\section{Related Work and Positioning}
\label{sec:related}

Chiplet integration combines smaller dies through advanced packaging to improve
reuse, yield options, and technology specialization, while introducing new
interconnect, test, thermal, and integration challenges
\cite{li2020chiplet,lau2021chiplets}. UCIe addresses the interoperability
problem with a layered die-to-die standard. Sharma \emph{et al.} describe the
architecture, circuit, channel, and packaging rationale behind UCIe 1.0
\cite{sharma2022ucie}; later work examines memory/storage applications
\cite{sharma2024memory} and circuit/package scaling toward three-dimensional
integration \cite{nature2024ucie3d}. Those studies establish UCIe's system and
physical context. This work does not reproduce their bandwidth, energy, or
channel claims; it targets the digital training-control evidence boundary.

Physical UCIe studies further illustrate why that boundary matters. Thermal and
signal-integrity analysis of UCIe-A channels \cite{son2023thermal} and
FPGA-based high-speed interface verification platforms \cite{chang2022fpga}
operate with channel, transceiver, or measurement evidence that is absent here.
Consequently, successful RTL simulation and FPGA fitting cannot imply BER,
electrical margin, or interoperability.

Verification practice commonly combines directed bring-up, constrained-random
exploration, scoreboards, assertions, and coverage. IEEE Std 1800.2 defines the
UVM framework \cite{ieee1800_2}. Vamsikrishna \emph{et al.} describe a unified
testbench supporting directed, random, and pseudo-random operation selection,
including a UCIe-oriented example \cite{vamsi2025unified}. Intel and Cadence
report PHY-model simulation interoperability in which sideband transactions and
high-level LTSM state stepping coordinate vector traffic
\cite{cadenceIntelUcie}. These public accounts demonstrate mature methodology
and commercial model use, but do not provide a directly comparable open RTL
controller, independent lane-LFSR/message-order predictor, retained-error CSR,
and fixed-seed evidence chain.

The contribution here is therefore the auditable combination rather than any
individual technique: a bounded UCIe 2.0-derived RTL control subset, an
interface-derived predictor and SVA layer, explicit negative scenarios, and
qualified FPGA evidence with a preserved failed-fit result. No “first” claim is
made; public literature and proprietary implementations continue to evolve,
including the newer UCIe 3.0 specification \cite{ucie30}.
